\section{Introduction}

An agent skill is a directory with a \texttt{SKILL.md} file, optional scripts and reference material that an LLM agent reads when a task calls for it \citep{anthropic2025introducingskills,agentskills2026spec}.\footnote{The authors used an LLM-based assistant (Claude, Anthropic) for code, analysis scripts and drafting. The authors checked every claim, number and citation, and take full responsibility for the content.} Public directories now index skills at large scale (one advertises more than 276{,}000 \citep{agentskill2026}), and registries have started to gate listing or signing on a with-and-without evaluation \citep{nvidia2026trustpipeline,skillmill2026}. Whether a skill helps is not obvious. SkillsBench found that curated skills raised the average pass rate from 33.9\% to 50.5\%, while skills an agent wrote for itself scored 8.1 to 11.5 points below no skill in the configurations tested \citep{li2026skillsbench}. SWE-Skills-Bench found no significant effect for 92 of 98 skills \citep{han2026sweskillsbench}, and WebDev-Skills-Bench found that injecting skills lowered mean Pass@2 by 1.3\% to 4.2\% \citep{yang2026webdevskills}.

A registry that admits skills is therefore running a hypothesis test, many times, on data that arrive sequentially, for authors who can resubmit. Each of these features defeats a familiar procedure used in the obvious way. A confidence interval re-checked as results arrive loses its coverage; a test applied to hundreds of skills without correction produces false discoveries at the nominal rate; and an author who resubmits a lightly edited skill until it passes will eventually succeed. Prespecified fixed-sample tests, multiplicity corrections and group-sequential designs remain valid when used as designed; the difficulty is that a registry needs all three protections at once, on exams that it wants to stop early. The registry's own record also has to be trustworthy. A reader should be able to check that the exam was not chosen after seeing the skill, that it was not stopped at a lucky moment, and that the reported decision follows from the transcripts.

The problem is not specific to skills. Model leaderboards, benchmark submission servers and any certification that lets a submitter retry face the same combination of sequential data, multiplicity and resubmission; skills are our motivating case and the source of our data. This paper asks two questions. First, can the per-skill evidence that current benchmarks collect support such decisions? Second, what does a decision procedure that is valid under optional stopping, multiplicity and resubmission cost, and how can a third party check it?

\paragraph{Findings.}
\begin{enumerate}
\item \textbf{Current per-skill evidence is too thin for sequential certification at library scale} (Section~\ref{sec:res-real}). In the released paired outcomes of SWE-Skills-Bench, \NSweGolden{} cells are ``golden'' under the source paper's uncorrected rule, \NSweBH{} survive Benjamini-Hochberg control, and no betting test admits any cell with 10 units. In SkillsBench, a one-off Benjamini-Hochberg analysis of fixed-sample $t$-tests certifies \NSbTBH{} of \NSbTasks{} skills, and they replicate on held-out configurations, so the signal is real. But a registry tests skills as they arrive. With 18 units, a per-exam betting test admits \NSbBet{} skills, but online control over skills in arrival order (e-LOND) keeps \NSbELond{} on average.
\item \textbf{Common decision rules do not hold their error rate} (Section~\ref{sec:res-validity}). A $t$-test re-run after each variant admits useless skills \NEOnePeekTMin\% to \NEOnePeekTMax\% of the time, and a 95\% interval checked every ten variants admits them \NEOnePeekCIMin\% to \NEOnePeekCIMax\% of the time, where each rule is meant to hold 5\% (the interval rule, being one-sided at 2.5\%, even less). A stale stored baseline breaks every rule, valid or not (\NEOneStaleMin\% to \NEOneStaleMax\%), so concurrent pairing matters more than the choice of test.
\item \textbf{Library-level control needs more than per-exam stopping followed by e-BH} (Section~\ref{sec:res-library}). Stopping each exam at $1/\alpha$ and then applying e-BH within a window admits nothing in our library simulations, because every stopped e-value is close to $1/\alpha$ and e-BH then needs most of the window to cross (Remark~\ref{prop:ebh}). A useless skill resubmitted \PResub{} times passes a per-exam gate \NResubPerExam\% of the time; alpha-spending over a skill's lineage holds this to \NResubLineage\%.
\item \textbf{Resubmission-proofness has a price set by an entropy} (Section~\ref{sec:theory}). One exam at level $a$ needs about $\log(1/a)/K_{\inf}$ hidden variants, where $K_{\inf}$ is the information per variant against the closest useless skill \citep{agrawal2025stopping}. We show that controlling error over a skill's resubmissions costs an honest author at least $H(\pi)/K_{\inf}$ extra variants, where $H(\pi)$ is the entropy of their attempt count, that spending according to $\pi$ attains this, and that a schedule $q$ different from $\pi$ adds $\mathrm{KL}(\pi\|q)/K_{\inf}$. The lower bound matches up to a constant, so choosing a schedule is a coding problem: we give the minimax schedule over a class of attempt-count laws and a schedule learned from earlier lineages, whose excess falls like $N/m$ after $m$ lineages and whose cumulative excess is within a factor of two of the best possible. The price is what the lineage guarantee costs on its own. In a registry that has run for 50 or more exams with online false-discovery control, the online level is the stricter one in all but \NLcBindMax\% of exams, so resubmission protection comes almost free and the schedule changes the cost by under \NLcSchedSpread\% (E19); the price binds for a new registry, for the first lineages, and for a registry that budgets lineages separately. A level-tuned mixture of bets keeps the price, with leading coefficient one and explicit constants at every finite level, when the gate does not know the skill's effect. On the released SkillsBench outcomes the bounds hold, and the standard capped plug-in bet needs a median \NThPrplRatio{} times (95\% bootstrap interval over skills \NThPrplRatioLo{} to \NThPrplRatioHi) the information lower bound $\log(1/\alpha)/K_{\inf}$ on the number of variants.
\item \textbf{Validity has a cost that can be planned} (Section~\ref{sec:res-cost}). Resampling the released SkillsBench outcomes, certifying the median helpful curated skill needs \NSbMedNAlpha{} hidden variants per exam, or \NSbMedNBonf{} with library-wide control (95\% bootstrap intervals over skills \NSbMedNAlphaLo--\NSbMedNAlphaHi{} and \NSbMedNBonfLo--\NSbMedNBonfHi), against the 18 paired units available today. These figures treat the benchmark's harness and model configurations as stand-ins for hidden variants, an exchangeability the released data cannot test, so they are planning estimates rather than guarantees.
\item \textbf{The gate works on live models and published skills} (Section~\ref{sec:res-live}). On generated hidden variants in two task families with Claude Haiku~5.5, it admitted both helpful skills, after \NLsHelpN{} and \NLsBHelpN{} variants, and refused placebos, a skill carrying answers to public examples and skills with wrong instructions, for under US\$7 in API fees, in families we designed so that the lifts were large. Repeating both studies with the open-weight Nemotron~3.5 Super VL gave the same decisions for US\$\NLnTotalCost. On three published skills by other authors, for APA references, AP news style and commit messages, it admitted every skill on both models, including the APA skill, whose lift was moderate, after \NLsFourRealN{} and \NLnFourRealN{} variants, for about US\$\NLFourTotalCost{} for the whole study; a placebo of general news-writing advice was also admitted, because it too raised the pass rate. On a task family built by Google DeepMind, three published math skills lowered both models' pass rates and the gate refused them, a placebo and a harmful skill, for US\$\NLFiveTotalCost{} in all; with Nemotron the harmful skill raised the pass rate slightly, by curbing runaway reasoning, and ran to the cap without admission. On the SkillsBench tasks, where harness failures left \NLoScored{} of \NLoTasks{} tasks scorable, it admitted nothing. Every decision can be recomputed from the ledger and transcripts, which will be released with the code.
\end{enumerate}

\paragraph{Contributions.} (i) A statement of skill admission as a sequential, multiple and resubmittable testing problem, with four requirements a registry decision should meet (Section~\ref{sec:setup}). (ii) A gate that meets them, assembled from existing tools: a paired exam on commit-reveal hidden variants, the betting test of \citet{waudbysmith2024betting} with a futility margin, e-LOND \citep{xu2024online} with lineage alpha-spending, and a dual-arm credential whose decision anyone can recompute from published transcripts (Section~\ref{sec:gate}). The statistical components are not new. Lineage spending is classical online family-wise control within one skill's resubmissions \citep{tian2021online}, and the failure of ``stop at $1/\alpha$, then e-BH'' for any bet cap is an elementary observation (Remark~\ref{prop:ebh}). What is new is their combination for admission under resubmission. (iii) A characterization of what resubmission-proofness costs an honest author, namely an entropy lower bound over lineage-valid gates whose exams are fixed in advance, attained by spending according to the law of attempt counts, with the mismatch cost of any other schedule, the minimax and learned schedules that follow from it, and an explicit finite-level version for a gate that does not know the effect (Section~\ref{sec:theory}). It builds on known bounds for the cost of one exam \citep{agrawal2025stopping,larsson2025numeraire}; the proof averages the one-exam bound over attempts, and what it adds is the explicit price of resubmission-proofness and the schedule that attains it. (iv) A reanalysis of released per-skill paired outcomes for admission decisions, a simulation study calibrated to them, and committed live studies on two models (Sections~\ref{sec:eval}--\ref{sec:results}). (v) An open reference implementation with a verifier and tests.

\paragraph{What we do not claim.} This work began as the design of a ``campus'' in which agents distill skills from shared traces and a skill enters a shared library only after an exam. Whether such peer-synthesized skills beat self-generated ones is not answered here (Appendix~\ref{app:campus} gives a study protocol). We also do not claim that per-skill testing or anytime-valid evaluation of LLMs is new: SWE-Skills-Bench already applies a per-skill significance test \citep{han2026sweskillsbench}, and anytime-valid procedures have been applied to self-evolving LLM systems and leaderboards \citep{yin2026enough,gao2026bbedge}.
